% from ega0-19-fr.tex
\oldpage[0\textsubscript{IV-1}]{94}

tout $n$ (Bourbaki, \emph{Alg. comm.}, chap.~III, §~2,
n\textsuperscript{o}~8, cor.~3 du th.~1)~; d'où $c)$ par passage à la limite
projective.
\end{env}

\begin{remark}[19.5.5]
\label{0.19.5.5-fr}
Sous les hypothèses générales de (19.5.4), supposons en outre que
$\mathfrak J$ soit un idéal \emph{maximal} de $B$, de sorte que
$C=B/\mathfrak J$ soit un corps, et que $\mathfrak J/\mathfrak J^2$ soit un
$C$-espace vectoriel de dimension \emph{finie}. Alors les conditions $a)$,
$b)$ et $c)$ de (19.5.4) sont encore équivalentes à la suivante~:
\begin{enumerate}
  \item[d)] \emph{Étant donnés une algèbre de polynômes}
  $E=C[T_1,\ldots,T_n]$ \emph{sur le corps} $C$, \emph{et en désignant par}
  $\mathfrak n$ \emph{l'idéal maximal de} $E$ \emph{engendré par les} $T_i$
  $(1\leq i\leq n)$, \emph{alors, pour tout idéal} $\mathfrak b$ \emph{de}
  $E$ \emph{contenant une puissance} $\mathfrak n^k$, \emph{tout
  homomorphisme} $v:B\to E/\mathfrak b$ \emph{de} $A$-\emph{algèbres}
  $C$-\emph{augmentées se factorise en}
  $B\xrightarrow{w}E/\mathfrak n^k\to E/\mathfrak b$, \emph{où} $w$ \emph{est
  un} $A$-\emph{homomorphisme}.
\end{enumerate}
En effet, $\mathfrak J/\mathfrak J^2$ étant ici un $C$-module libre, il suffit
de vérifier que la condition $d)$ entraîne que l'homomorphisme canonique
$\varphi:\mathbf S_C^*(\mathfrak J/\mathfrak J^2)\to\gr_{\mathfrak J}^*(B)$
est bijectif. Or ici $\mathbf S_C^*(\mathfrak J/\mathfrak J^2)$ s'identifie à
l'algèbre de polynômes $E=C[T_1,\ldots,T_n]$, où
$n=\operatorname{rg}_C(\mathfrak J/\mathfrak J^2)$, et l'idéal d'augmentation
de $E$ est l'idéal $\mathfrak n$ engendré par les $T_i$. Pour tout entier
$k\geq0$, l'hypothèse que $C$ est une $A$-algèbre discrète formellement lisse
entraîne, comme dans (19.5.4.2), l'existence d'un $A$-homomorphisme surjectif
$f:E/\mathfrak n^{k+1}\to B/\mathfrak J^{k+1}$, tel que
$\gr^j(f)=\varphi_j$ pour tout $j\leq k$. Si
$\mathfrak b/\mathfrak n^{k+1}=\operatorname{Ker}(f)$,
$B/\mathfrak J^{k+1}$ s'identifie donc à $E/\mathfrak b$, et l'hypothèse $d)$
permet de factoriser l'homomorphisme canonique $B\to B/\mathfrak J^{k+1}$ en
$B\xrightarrow{w}E/\mathfrak n^{k+1}\xrightarrow{f}B/\mathfrak J^{k+1}$~;
comme $\gr^0(f)$ est l'identité, on a
$w(\mathfrak J)\subset\mathfrak n/\mathfrak n^{k+1}$, donc
$w(\mathfrak J^{k+1})=0$~; on conclut comme dans (19.5.4.2) que $\varphi_k$
est injectif pour tout $k$, ce qui achève de prouver notre assertion.
\end{remark}

\begin{env}[19.5.6]
\label{0.19.5.6-fr}
\emph{Démonstration du théorème} (19.5.3). --- Soit
$(\mathfrak b_\alpha)$ un système fondamental décroissant d'idéaux ouverts de
$B$. Nous poserons
\[
  B_\alpha=B/\mathfrak b_\alpha,\qquad
  C_\alpha=B/(\mathfrak b_\alpha+\mathfrak J),\qquad
  \mathfrak J_\alpha=(\mathfrak b_\alpha+\mathfrak J)/\mathfrak b_\alpha,
\]
de sorte que
\[
  C_\alpha=B_\alpha/\mathfrak J_\alpha\qquad\text{et}\qquad
  \mathfrak J_\alpha^n/\mathfrak J_\alpha^{n+1}
  =(\mathfrak b_\alpha+\mathfrak J^n)/(\mathfrak b_\alpha+\mathfrak J^{n+1}).
\]
Les $C$-modules
$((\mathfrak b_\alpha\cap\mathfrak J^n)+\mathfrak J^{n+1})/
\mathfrak J^{n+1}$ de $\mathfrak J^n/\mathfrak J^{n+1}$ forment un système
fondamental de sous-modules ouverts du $C$-module topologique
$\mathfrak J^n/\mathfrak J^{n+1}$~; comme
$(\mathfrak b_\alpha\cap\mathfrak J^n)+\mathfrak J^{n+1}
=\mathfrak J^n\cap(\mathfrak b_\alpha+\mathfrak J^{n+1})$, on a
\[
  \mathfrak J^n/((\mathfrak b_\alpha\cap\mathfrak J^n)+\mathfrak J^{n+1})
  =\mathfrak J^n/(\mathfrak J^n\cap(\mathfrak b_\alpha+\mathfrak J^{n+1}))
  =(\mathfrak b_\alpha+\mathfrak J^n)/(\mathfrak b_\alpha+\mathfrak J^{n+1})
  =\mathfrak J_\alpha^n/\mathfrak J_\alpha^{n+1}.
\]
\end{env}

\begin{env}[19.5.6.1]
\label{0.19.5.6.1-fr}
Soient $P$, $Q$ deux $C_\alpha$-modules discrets,
$u:P\to Q$ un $C_\alpha$-homomorphisme surjectif. L'anneau discret
$B/(\mathfrak b_\alpha+\mathfrak J^2)$ est une $B$-extension de $C_\alpha$ par
l'idéal de carré nul $\mathfrak J_\alpha/\mathfrak J_\alpha^2$. Soit
$v:\mathfrak J/\mathfrak J^2\to Q$ un $C$-homomorphisme continu~; remplaçant
au besoin $\mathfrak b_\alpha$ par un idéal ouvert plus petit de $B$, on peut
supposer que le noyau de $v$ contient le sous-$C$-module ouvert
$((\mathfrak b_\alpha\cap\mathfrak J)+\mathfrak J^2)/\mathfrak J^2$~;
passant au quotient, on en déduit un $C_\alpha$-homomorphisme de modules
discrets $v':\mathfrak J_\alpha/\mathfrak J_\alpha^2\to Q$~; soit
$E_\alpha$ la $B_\alpha$-extension de $C_\alpha$ par $Q$ déduite de
$B/(\mathfrak b_\alpha+\mathfrak J^2)$ au moyen de $v'$, et soit
$v'_*:B/(\mathfrak b_\alpha+\mathfrak J^2)\to E_\alpha$ le
$B_\alpha$-homomorphisme correspondant (18.2.8)~; $E_\alpha$ est une
$A$-algèbre topologique discrète, et l'isomorphisme canonique
$C_\alpha\xrightarrow{\sim}E_\alpha/Q$ donne par composition un
$A$-homomorphisme continu $f:C\to C_\alpha\to E_\alpha/Q$. Mais comme $Q$ est
de carré nul dans $E_\alpha$ et que $C$ est une $A$-algèbre formellement
lisse, $f$ se factorise en $C\xrightarrow{g}E_\alpha\to E_\alpha/Q$ où $g$
est un $A$-homomorphisme continu. Comme $g$ est continu et $E_\alpha$ discret,
il existe un $\beta\geq\alpha$ tel que $g$ se factorise en
$C\to C_\beta\xrightarrow{g'}E_\alpha$. D'autre part, soit $E_\beta$ la
$A$-extension de $C_\beta$ par $Q$, image réciproque de $E_\alpha$ par
l'homomorphisme canonique $C_\beta\to C_\alpha$~; l'existence de $g'$ (tel que
le composé $C_\beta\to E_\alpha\to C_\alpha$ soit l'homomorphisme canonique)
équivaut au fait que $E_\beta$ est une extension $A$-\emph{triviale}, et peut
donc s'identifier à $\mathrm D_{C_\beta}(Q)$. Cela étant, l'homomorphisme
surjectif $u:P\to Q$ définit canoniquement un homomorphisme surjectif
$\mathrm D_{C_\beta}(P)\to\mathrm D_{C_\beta}(Q)$ (18.2.9), dont le noyau est
un idéal $\mathfrak L$ de l'extension $E'_\beta=\mathrm D_{C_\beta}(P)$
contenu dans $P$, et \emph{a fortiori} de carré nul. Soit
$f':B\to E_\beta=E'_\beta/\mathfrak L$ le $A$-homomorphisme continu définissant
la structure de $B$-algèbre topologique de $E_\beta$~; comme $\mathfrak L$ est
de carré nul et que $B$ est une $A$-algèbre formellement lisse, $f'$ se
factorise en $B\xrightarrow{h'}E'_\beta\to E'_\beta/\mathfrak L$, où $h'$ est
un $A$-homomorphisme continu. La construction de $E'_\beta$ donne par
composition un $A$-homomorphisme $\psi:E'_\beta\to E_\beta\to E_\alpha$, et
il est clair que le diagramme
\[
  \xymatrix{
    E'_\beta \ar[r]^{\psi} & E_\alpha \\
    B \ar[u]^{h'} \ar[r]_{\theta} &
    B/(\mathfrak b_\alpha+\mathfrak J^2) \ar[u]_{v'_*}
  }
\]
est commutatif. D'ailleurs, l'image par $\psi\circ h'$ de $\mathfrak J$ est
contenue dans $Q$, donc l'image par $h'$ de $\mathfrak J$ est contenue dans
$P$, et l'image de $\mathfrak J^2$ par $h'$ est nulle. Restreignant $h'$ et
$\theta$ à $\mathfrak J$, on obtient un diagramme commutatif
\[
  \xymatrix{
    P \ar[r]^{u} & Q \\
    \mathfrak J/\mathfrak J^2 \ar[u]^{w} \ar[ur]_{v}
  }
\]
où $w$ est continu, ce qui prouve que
$\mathfrak J/\mathfrak J^2$ est un $C$-module formellement projectif.
\end{env}

